% paper/section_7_limitations.tex
% Limitations. Drawn from PAPER_NOTES_JUNE3 + the June 10 audit caveats.
% Input under a \section{Limitations} declared in the skeleton.

\paragraph{No head-to-head with ClawGuard specifically.} An earlier version
of this work lacked any AgentDojo evaluation at scale; that general gap is
now closed by the 629-real-trajectory validation in \S\ref{sec:gpt4o}. A
narrower gap remains: we do not run a controlled head-to-head against
ClawGuard's implementation (\S\ref{sec:related}), which intercepts live
sessions with no dataset-replay harness compatible with our pipeline.

\paragraph{External tool-call-content baselines are structurally
inapplicable, not weaker.}
\label{par:external-baselines}
We additionally attempted comparison against the two closest available
external tools at the tool-call content level. Both proved structurally
inapplicable rather than merely weaker: CodeShield's supported-language set
(C/C++/C\#/Java/JavaScript/PHP/Python/Rust) excludes the shell, PowerShell,
and cmd dialects comprising $\sim$87\% of our reference corpus's
adversary-emulation commands, and gitleaks --- a secret-value scanner ---
produces zero findings on all 519 cases, since they encode
credential-access techniques rather than embedded literal secrets (its one
finding across 848 benign calls was an empty-string hash misread as an API
key). We report this as corroborating the enforcement-boundary gap this
paper targets, not as a performance comparison.

\paragraph{Tamper-evident, not tamper-proof.} The hash-chained ledger detects
post-hoc edits to any recorded decision or its provenance, but an attacker
with write access to the database can recompute the entire chain; defeating
that requires anchoring the chain head to an external append-only medium,
which we do not implement. We found and fixed a real bug in the ledger's
own \emph{verification} query during this work (\S\ref{sec:provenance}) ---
the kind of bug that only surfaces when the mechanism is run end to end
against real data, not when merely reviewed as a design.

\paragraph{Curated corpus and deterministic rule set.} Both the cosine-triage
corpus and the deterministic prefilters' rule sets are authored against
published taxonomies (MITRE ATT\&CK/ATLAS, AtomicRedTeam, GTFOBins) and known
attack corpora (InjecAgent, AgentDojo). A genuinely novel technique outside
both may not be caught; \S\ref{sec:generalization-audit} characterizes this
precisely for each minister rather than leaving it as an untested assumption.

\paragraph{The reserve\_hotel gap is representative of a broader,
unfixed class.} \S\ref{sec:gpt4o} categorizes the 19 genuine real-trajectory
misses in full: five (\texttt{reserve\_hotel}) are a scoped CHANNEL
destination-tool gap whose natural fix we tried and reverted, since it
introduces a new false-positive class; the remaining fourteen have no
destination value to trace at all and are structurally not closeable by any
refinement of CHANNEL's mechanism --- exactly the shape of question
\S\ref{par:navigator-authflow} below reports as unresolved.

\paragraph{Single embedding model.} The cosine-triage layer and NAVIGATOR use
a single encoder (\texttt{bge-base-en-v1.5}), selected for open-source
offline reproducibility rather than by a comparative ablation; whether an
alternative open-source retrieval encoder would shift the residual
detection rate is an open question we leave to future work.

\paragraph{PowerShell's real AST is correct but too slow for the live path.}
\S\ref{sec:vault-executor} reports a measured $\sim$1.4\,s per call for a
real PowerShell grammar parse via subprocess, 18$\times$ over budget, so
PowerShell calls are scored by the structural tokenizer rather than a real
AST live; the real parser remains available as an offline/audit tool. A
faster in-process PowerShell AST is a scoped next step we did not pursue
this cycle.

\paragraph{NAVIGATOR authorization layer: a hard-constraint layer still in
progress, after three failed heuristics.}
\label{par:navigator-authflow}
NAVIGATOR is the one minister whose deterministic rearchitecture is not yet
complete. A direct deterministic swap mirroring VAULT/EXECUTOR/CHANNEL
showed zero net discrimination gain and was abandoned; a second attempt
reframes the question as authorization rather than similarity (a pinned,
immutable session intent, classifying each tool as observational or
consequential, then a hard-constraint layer deciding whether a
consequential action is \emph{derivable} from that intent). The first two
pieces --- pinned-intent capture and the taxonomy --- are complete and
validated, including under the generalization audit below. The
hard-constraint layer failed three independently designed heuristics
against a frozen discriminability criterion: (1)~tool-name/instruction-text
overlap, evaded by attacker text sharing surface words with the legitimate
request;
(2)~argument-value overlap alone, wrongly flagging legitimate values read
earlier in the session; (3)~a provenance-aware version of (2), composed
with CHANNEL's own session state, still over the frozen false-positive
threshold. Reported precisely because a negative result with a specific
failure mechanism is more useful to future work than a vague one.

\paragraph{A proposed alternative for NAVIGATOR: intent-drift, explored and
not validated.} Independently of the hard-constraint layer above, we
explored a second candidate mechanism for NAVIGATOR --- flagging a tool call
as drift when its embedding similarity to the session's original
user-stated goal falls below a threshold --- which reached 95.2\% recall at
0\% false-positive rate on an initial 100-case dataset (83 attack, 17
benign), but did not generalize: enlarging the benign sample from 17 to 39
with real multi-step AgentDojo sessions raised the false-positive rate to
35.9\% at the same threshold, because ordinary mid-session follow-up actions
(e.g., checking a balance right after a payment) are semantically distant
from the session's \emph{original} goal despite being legitimate, and
held-out recall on the enlarged set collapsed to 22--56\%. This is
consistent with an independent evaluation of the same mechanism on
AgentDojo ground-truth, which measured 7.3\% recall. We report this as
future work, not a contribution: NAVIGATOR's deployed detector remains the
cosine-corpus mechanism in \S\ref{sec:generalization-audit}, unchanged.

\paragraph{No adaptive-attacker evaluation, and a first-pass obfuscation
check.} Every attack corpus in this paper is static: no attacker knows
Kavach's rules or adapts to them. Deterministic matching is historically
evadable --- encoding, string construction, variable indirection --- and
while bashlex's AST closes a class of these for sh/bash, PowerShell and
cmd remain on a tokenizer. First-pass check: 335 previously-detected
sh/bash cases through three transforms (base64 encode-and-pipe,
variable-split, variable indirection). Bottom line held at 100\% (0/335
missed). After correctly distinguishing genuine evasion from TRAJECTORY's
session-level re-confirmation merely relabeling an already-correct verdict
(\S\ref{sec:design-tradeoffs}), true per-mechanism evasion is 7.0\% (encode,
13/185), 4.0\% (split, 3/75), 25.3\% (indirect, 19/75); every evasion was
caught by a different mechanism. One backup is itself imprecise: encode's
backstop is a single NAVIGATOR pattern matching the literal
\texttt{echo~BASE64~|~base64~-d~\allowbreak|~sh} shape --- confirmed
directly to also block an equivalent benign command, so it is a blanket
rule, not intent-aware detection. 335 cases across three hand-written
transforms is a first pass, not adversarial; an attacker targeting two or
more mechanisms' blind spots at once is untested. Until that broader study
is done, \S\ref{sec:eval}'s rates should be read as non-adaptive.

\subsection{Generalization beyond the benchmark vocabulary}
\label{sec:generalization-audit}

Kavach's deterministic rules were, in large part, authored by mining
detection misses on the exact benchmarks evaluated against
(\S\ref{sec:kavach-pb}), creating a real risk of circular validation. We
therefore ran an independent adversarial audit before submission against a
fictional agent environment (``AtriumOS'') with an entirely invented tool
vocabulary sharing no tool name, corpus term, or surface phrasing with
InjecAgent, AgentDojo, or our reference corpus: 18 attack and 18
matched-benign novel-vocabulary scenarios through the real detection
pipeline. Summary verdict: EXECUTOR generalizes cleanly, VAULT mostly does
(one architecture-inherent limit), CHANNEL's escalation path generalizes
but its hard-block path does not, and NAVIGATOR's cosine mechanism does not
transfer to novel vocabulary at all. Full methodology and per-minister
results in Appendix~A.
